\section{Related Literature}
\label{sec:literature}

The paper relates to three strands of work. First, a growing literature studies strategic industrial policy across jurisdictions. \citet{ChenLi2026} develop a quantitative spatial-equilibrium model of local industrial subsidy competition and compare local, cooperative, and central policy under varying market integration. \citet{WangXuYangZhu2024} solve non-cooperative tariffs and sector-specific industrial subsidies across many economies and sectors. These contributions establish that industrial policy can be strategic and multi-sectoral. Our object is narrower: each government has a fixed total policy capacity and chooses its composition across activities. The distortion therefore survives even when the policy total is held constant across decentralized and coordinated regimes.

Second, work on green and innovation policy studies upstream and downstream interventions and cross-border policy spillovers. \citet{Fischer2017} examines strategic green industrial policy when governments value domestic production, while \citet{FischerGreakerRosendahl2018} distinguish upstream technology support from downstream deployment support in a two-region setting. \citet{BorotaMilicevicEtAl2026} analyze cooperation in regional R\&D subsidies with cross-border knowledge diffusion. Our result does not rest on the novelty of upstream/downstream policy or coordination itself. The distinctive object is a fixed-capacity portfolio choice that generates a strict difference between the decentralized and coordinated thresholds for reallocating policy toward a complementary activity.

Third, the paper connects to place-based policy and smart specialization. Smart-specialization strategies explicitly require regions to concentrate policy on a limited set of priorities and related transformative activities \citep{ForayEichlerKeller2021}. Recent evidence from China shows that local governments systematically favor industries aligned with their own specialization and do not display the same favoritism toward supply-chain-linked industries or specialization extending across jurisdictional boundaries \citep{TianWangZhang2027}. These observations motivate the model's combination of scarce policy capacity, a local return premium, and cross-jurisdiction complementarity. They do not identify the model parameter values or establish the welfare theorem; those remain theoretical objects.

The contribution is therefore not that local governments compete, that industrial policy can duplicate activity, or that coordination can internalize spillovers. It is the composition-specific mechanism: incomplete local capture of complementary cross-regional value can make decentralized governments retain the same locally attractive priority after a coordinated reallocation of the same fixed policy capacity has become welfare improving.
